% Paper 30: SU(2) Wilson Lattice Gauge Structure on the GeoVac Hopf Graph
% GeoVac Project, April 2026
% v1.0 -- Non-abelian sibling of Paper 25.

\documentclass[11pt]{article}
\usepackage{amsmath,amssymb,amsthm}
\usepackage{booktabs}
\usepackage{array}
\usepackage[margin=1in]{geometry}

\newtheorem{theorem}{Theorem}
\newtheorem{proposition}{Proposition}
\newtheorem{corollary}{Corollary}
\newtheorem{observation}{Observation}
\newtheorem{definition}{Definition}
\newtheorem{remark}{Remark}

\title{SU(2) Wilson Lattice Gauge Structure on the GeoVac Hopf Graph:\\
The Non-Abelian Sibling of Paper 25}
\author{GeoVac Project}
\date{April 2026}

\begin{document}
\maketitle

\begin{abstract}
Paper~25~\cite{paper25} placed the GeoVac Hopf graph inside the
Wilson--Hodge lattice-gauge vocabulary with a $U(1)$ connection on
edges, and explicitly left as open (\S VII.A) the question of whether
a non-abelian gauge structure is natural on the same graph.  This paper
answers that question for the compact simply-connected non-abelian Lie
group of rank~1, $\mathrm{SU}(2)$.  The coincidence $S^3 =
\mathrm{SU}(2)$ (as manifolds) makes this case structurally special:
the underlying manifold of the GeoVac graph \emph{is} the gauge group.
Taking link variables $U_e \in \mathrm{SU}(2)$, plaquettes as the
primitive closed non-backtracking walks of the Ihara zeta
setup~\cite{paper29}, and the Wilson action
$S_W[U] = \beta \sum_P \bigl(1 - \tfrac{1}{2}\mathrm{Re}\,\mathrm{tr}\,
U_P\bigr)$, we obtain a gauge-invariant finite-dimensional lattice
gauge theory with Haar-normalizable partition function.
Three structural results follow.
(1)~\emph{$U(1)$ reduction to Paper 25.}  The diagonal (maximal-torus)
sector of $\mathrm{SU}(2)$ link variables reproduces Paper~25's
$U(1)$ Wilson--Hodge action to machine precision.  Paper~25 is the
Cartan-subalgebra projection of the present theory; the abelian
construction it documents is a structural sub-sector of a genuinely
non-abelian one.
(2)~\emph{The weak-coupling kinetic term completes the Hodge-1
Laplacian.}  Linearizing $S_W$ around the trivial vacuum $U_e = I$
produces, at second order in the gauge field, the plaquette (curl)
quadratic form $(\beta/8) \sum_a \langle A^a, B_2 B_2^{\top} A^a\rangle$
on $\mathfrak{su}(2)$-valued $1$-cochains, with per-plaquette
per-component coefficient $1/8 = 1/(4 N_c)$.  This curl form is
\emph{not} a scalar multiple of Paper~25's edge Laplacian
$L_1 = B^{\top} B$ (the two have complementary supports: the plaquette
columns of $B_2$ span $\ker L_1$ exactly); it is the co-exact part that completes
$L_1$ to the full discrete Hodge-1 Laplacian
$B^{\top} B + B_2 B_2^{\top}$ of the 2-complex with the Ihara
plaquettes attached as faces, exactly the face structure
Paper~25~\S II anticipates.
(3)~\emph{Plaquette enumeration and Monte Carlo Wilson-loop data.}
On the $S^3$ Coulomb graph at $n_{\max} = 2, 3, 4$, the first-Betti
counts are $\beta_1 = 0, 2, 8$ and the primitive 4-cycle counts are
$0, 2, 8$; a 2000-sample Metropolis run on $n_{\max} = 3$ at $\beta
\in \{0.5, 1, 2, 5\}$ gives Wilson expectations $\langle W \rangle =
0.126, 0.158, 0.332, 0.681$, monotonic with $\beta$ and
matched to the leading-order character expansion at the $2$--$35\%$
level across the four $\beta$ values (closest at the strong- and
weak-coupling ends).  All 36 tests in \texttt{tests/test\_su2\_wilson\_gauge.py}
pass.
The scope is a compact-graph non-abelian lattice gauge construction,
not a Yang--Mills mass-gap claim on $\mathbb{R}^4$; the paper closes
with a synthesis observation that the GeoVac framework now carries
\emph{five} action principles---maximum entropy on nodes, Rayleigh--Ritz
on node-wavefunctions, Wilson action on gauge edges, Ihara path-integral
on walks, and spectral action on Dirac spectrum---suggesting---though we do not claim---a framework-wide
organizing principle of least action.
An $\mathrm{SU}(3)$ sibling on the Bargmann--Segal $S^5$ graph
(Paper~24) has been constructed in Sprint~ST-SU3
(\texttt{debug/st\_su3\_wilson\_memo.md}, May~2026), confirming
Proposition~1's analog (Cartan reduction to $U(1)\times U(1)$) and
Proposition~2's analog (kinetic coefficient $1/12 = 1/(4 N_c)$) at machine
precision; matter coupling on $S^5$ encounters the Sprint~5 Track~S5
Clebsch--Gordan-intertwiner obstruction, sharpening Paper~24's
Coulomb/HO asymmetry to four layers (spectrum-computing $L_0$,
calibration $\pi$, Wilson gauge with natural matter coupling, and the
universal Wilson--Hodge vocabulary; the first three Coulomb-specific,
the fourth universal).
\end{abstract}

% ===================================================================
\section{Introduction}
\label{sec:intro}
% ===================================================================

The GeoVac framework discretizes Coulomb bound-state quantum mechanics
onto a finite graph that converges to the unit three-sphere $S^3$ in
the continuum limit~\cite{paper7}.  Paper~25~\cite{paper25} observed
that the same graph, read through its signed node--edge incidence
matrix $B$, carries all the combinatorial ingredients of a lattice
gauge theory: nodes as matter sites, edges as link variables in a
$U(1)$ connection, node Laplacian $L_0 = B B^{\top}$ as matter
propagator, edge Laplacian $L_1 = B^{\top} B$ as discrete Hodge-1
Laplacian.  Section~VII.A of Paper~25 explicitly flagged as open the
question of whether a non-abelian gauge structure is also natural on
the same graph, and sharpened the question by an analysis of the
$S^5$ Bargmann--Segal graph of Paper~24~\cite{paper24} which
\emph{fails} to admit an $\mathrm{SU}(3)$ lattice gauge analog---the
$(N, 0)$ representation tower has Clebsch--Gordan intertwiners on
links, not fixed group elements, so it fails the Wilson
fixed-group-on-every-link requirement.

The $S^3$ case is categorically different.  As manifolds,
$S^3 = \mathrm{SU}(2)$: the underlying manifold of the GeoVac graph
\emph{is} the gauge group.  Taking $\mathrm{SU}(2)$-valued link
variables is therefore not a choice made by matching representation
content across shells; it is the natural object carried by the
manifold itself.  This is the cleanest geometric setting in which a
non-abelian Wilson lattice gauge theory can be built on a GeoVac graph.

This paper reports the construction.  The mathematics is
standard~\cite{wilson1974,creutz1983}\footnote{Concurrent work in
the Bonn/DESY Hamiltonian-lattice program develops SU(2) Hamiltonian
lattice gauge theory using discrete partitionings of the gauge group,
with motivation toward quantum simulation: see
Garofalo et al.~\cite{garofalo2023_su2_s3} for the SU(2) Hamiltonian
formulation based on $S_3$ partitionings (here $S_3$ denotes the
symmetric group, distinct from the manifold $S^3$ used here), and
Jakobs et al.~\cite{jakobs2025_su2_dynamics} for dynamical
continuum-limit studies.}: link variables
$U_e \in \mathrm{SU}(2)$, plaquette holonomies $U_P = \prod_{e \in P}
U_e$, Wilson action
$S_W = \beta \sum_P (1 - \tfrac{1}{2} \mathrm{Re}\,\mathrm{tr}\,U_P)$,
node-local gauge transformations $U_e \mapsto g_{s(e)}^{\phantom{\dagger}}
U_e\, g_{t(e)}^{\dagger}$, and Haar measure on each link.  What is new
is (i)~the choice of plaquettes as the primitive closed non-backtracking
walks of Paper~29~\cite{paper29}'s Ihara zeta machinery, rather than the
elementary four-cycles of a hypercubic lattice, (ii)~the confirmation
that the maximal-torus reduction is \emph{exactly} Paper~25's $U(1)$
construction, and (iii)~the identification of the weak-coupling kinetic term as the
plaquette (curl) quadratic form that completes Paper~25's edge
Laplacian $L_1$ to the full discrete Hodge-1 Laplacian of the
plaquette-filled 2-complex.

\paragraph*{Scope.}  This is an observation paper.  What it establishes
is a self-consistent non-abelian lattice gauge structure on the finite
compact GeoVac Hopf graph, together with three structural results.
What it does \emph{not} establish is a continuum-limit Yang--Mills mass
gap on $\mathbb{R}^4$, confinement (the graph has a single plaquette
class at the sizes tested, making the area/perimeter-law diagnostic
inapplicable), or any new derivation of $\alpha$ in the Paper~2 sense.
The rhetoric follows the CLAUDE.md~\S1.5 rule: the mathematical
construction is sharp; the physical interpretation is modest.  The
construction is a non-abelian instance of the Marcolli--van Suijlekom
gauge-network framework~\cite{marcolli_vs2014}, which defines finite
spectral triples on graphs whose spectral action is a Wilson lattice
gauge theory; Perez-Sanchez~\cite{perez_sanchez2025}
establishes that the continuum limit in that lineage is Yang--Mills
without Higgs, consistent with the scope statements of
Section~\ref{sec:scope}.

\paragraph*{Organization.}  Section~\ref{sec:recap} recaps the
Paper~25 Wilson--Hodge vocabulary in the form needed.
Section~\ref{sec:su2} states the $\mathrm{SU}(2)$ link variables,
plaquette structure, Wilson action, and gauge transformation.
Section~\ref{sec:characters} gives the character expansion of the
partition function on $\mathrm{SU}(2)$.  Section~\ref{sec:results}
reports the three structural results: $U(1)$-reduction exactness,
the kinetic-term / Hodge-completion result, and plaquette/Wilson-loop
data.
Section~\ref{sec:action_principles} places the construction inside a
framework-wide synthesis: the GeoVac project now exhibits five
independent action principles, each organizing a different sector of
its data.  Section~\ref{sec:scope} states explicitly what is and is
not claimed (Yang--Mills adjacency).  Section~\ref{sec:open} lists
open questions, and Section~\ref{sec:conclusion} concludes.

% ===================================================================
\section{Recap: The GeoVac Hopf Graph as a Lattice Gauge Structure}
\label{sec:recap}
% ===================================================================

Paper~25~\cite{paper25} specifies the combinatorial data in the form
required here.  We collect the definitions and fix notation; no new
claim is made in this section.

The GeoVac graph at cutoff $n_{\max}$ has vertex set $V$ indexed by
quantum-number triples $(n, l, m_l)$ with $n = 1, \ldots, n_{\max}$,
$l = 0, \ldots, n-1$, $m_l = -l, \ldots, +l$, and $|V| =
\sum_{n=1}^{n_{\max}} n^2$.  Edges are of two kinds: $L_{\pm}$
angular ladders connecting $(n, l, m_l) \leftrightarrow
(n, l, m_l \pm 1)$ at fixed $(n, l)$, and $T_{\pm}$ radial ladders
connecting $(n, l, m_l) \leftrightarrow (n', l, m_l)$ for
$n' = n \pm 1$ (with selection-rule variants~\cite{paper7,paper14}).
Fix an orientation on each undirected edge; the signed incidence
matrix $B \in \mathbb{Z}^{|V| \times |E|}$ has $B_{v, e} = +1$ if edge
$e$ starts at $v$, $-1$ if it ends at $v$, and $0$ otherwise.  The
combinatorial Hodge Laplacians~\cite{lim2020,schaub2020} are
\begin{equation}
L_0 = B B^{\top}, \qquad L_1 = B^{\top} B,
\label{eq:hodge_laplacians}
\end{equation}
the node and edge Laplacians respectively.  The first Betti number
\begin{equation}
\beta_1 = |E| - |V| + c
\label{eq:betti1}
\end{equation}
($c$ = number of connected components) equals the kernel dimension of
$L_1$ and counts the independent cycle classes.

The orienting convention for plaquettes differs here from the ``elementary
4-cycle'' convention of a hypercubic lattice.  We take plaquettes to be
the \emph{primitive closed non-backtracking walks} of the graph, i.e.\
the Ihara cycles of Paper~29~\cite{paper29}.  A primitive closed
non-backtracking walk on $L$ edges is a sequence $e_1 e_2 \cdots e_L$
with $\mathrm{head}(e_i) = \mathrm{tail}(e_{i+1})$, no $e_{i+1}$ equal
to the reverse of $e_i$ (including across the closure $e_L \to e_1$),
and not a cyclic rotation of a shorter closed walk.  This is the
natural Wilson generalization on a non-hypercubic graph: primitive
closed non-backtracking walks are exactly the irreducible ``boundary
of a disk'' objects whose holonomies enter the Wilson action.

Table~\ref{tab:plaquettes} collects the plaquette counts on the
$S^3$ Coulomb graph at $n_{\max} = 2, 3, 4$.  At $n_{\max} = 2$
the graph is a forest ($\beta_1 = 0$), so there are no plaquettes
and the $\mathrm{SU}(2)$ Wilson theory is trivial ($Z(\beta) = 1$).
At $n_{\max} = 3$ the theory is non-trivial for the first time: two
primitive 4-cycles in the $p$-block, one primitive 6-cycle, one
primitive 8-cycle.  At $n_{\max} = 4$ the graph has
$\beta_1 = 8$ and a richer plaquette structure.

\begin{table}[t]
\centering
\caption{Plaquette enumeration on the $S^3$ Coulomb graph.  $V, E, c$
are vertex, edge, and component counts; $\beta_1 = E - V + c$ is the
first Betti number; $N_L$ is the number of primitive closed
non-backtracking walks of length $L$ (one representative per
unoriented class).}
\label{tab:plaquettes}
\begin{tabular}{rrrrrrrr}
\toprule
$n_{\max}$ & $V$ & $E$ & $c$ & $\beta_1$ & $N_4$ & $N_6$ & $N_8$ \\
\midrule
2 & 5 & 3 & 2 & 0 & 0 & 0 & 0 \\
3 & 14 & 13 & 3 & 2 & 2 & 1 & 1 \\
4 & 30 & 34 & 4 & 8 & 8 & 7 & 18 \\
\bottomrule
\end{tabular}
\end{table}

The two primitive 4-cycles at $n_{\max} = 3$ both live in the
$p$-block.  With the canonical representative of each cycle chosen to
begin at its lexicographically smallest oriented edge, they are
\begin{align}
P_A & : (2, 1, -1) \to (2, 1, 0) \to (3, 1, 0) \to (3, 1, -1) \to (2, 1, -1),
\nonumber\\
P_B & : (2, 1, 0) \to (2, 1, 1) \to (3, 1, 1) \to (3, 1, 0) \to (2, 1, 0).
\label{eq:p_block_plaquettes}
\end{align}
These share the undirected edge $\{(2, 1, 0), (3, 1, 0)\}$.  The
plaquette-sharing structure becomes important in
Section~\ref{sec:characters}: the leading-order character expansion
treats plaquettes as independent, but the shared edge introduces
higher-order corrections.

% ===================================================================
\section{SU(2) Link Variables and the Wilson Action}
\label{sec:su2}
% ===================================================================

\subsection{Link variables}

Each undirected edge $\{u, v\}$ with chosen orientation $u \to v$
carries a link variable $U_{u \to v} \in \mathrm{SU}(2)$.  The reverse
edge carries
\begin{equation}
U_{v \to u} = U_{u \to v}^{\dagger},
\label{eq:reverse_convention}
\end{equation}
so that each undirected edge stores one independent $\mathrm{SU}(2)$
element.  An $\mathrm{SU}(2)$ matrix is parametrized by a unit
quaternion $(a, b, c, d)$ with $a^2 + b^2 + c^2 + d^2 = 1$ as
\begin{equation}
U = \begin{pmatrix} a + i d & c + i b \\ -c + i b & a - i d \end{pmatrix},
\qquad
\det U = 1, \quad U U^{\dagger} = I.
\label{eq:su2_quaternion}
\end{equation}
Equivalently, $U = \exp\bigl(i \theta\,\hat{n} \cdot \boldsymbol{\sigma}\bigr)$
where $\boldsymbol{\sigma} = (\sigma_x, \sigma_y, \sigma_z)$ are the
Pauli matrices, $\hat{n}$ is a unit 3-vector, and $\theta$ is the
half-rotation angle.

\subsection{Plaquette holonomy}

For a plaquette $P$ traversed in the ordered sequence
$e_1, e_2, \ldots, e_L$ (each $e_i$ an oriented edge), the plaquette
holonomy is the ordered product
\begin{equation}
U_P = U_{e_1} U_{e_2} \cdots U_{e_L}.
\label{eq:plaquette_holonomy}
\end{equation}
Since $\mathrm{SU}(2)$ is closed under matrix multiplication,
$U_P \in \mathrm{SU}(2)$.  The plaquette character
\begin{equation}
\chi(U_P) = \tfrac{1}{2} \mathrm{Re}\,\mathrm{tr}\,U_P
= \cos \theta_P,
\label{eq:plaquette_character}
\end{equation}
where $\theta_P$ is the half-angle of $U_P$ in its $\mathrm{SU}(2)$
rotation parametrization, takes values in $[-1, +1]$ with
$\chi(U_P) = +1$ iff $U_P = I$.

\subsection{Wilson action}

The Wilson action is
\begin{equation}
S_W[U] = \beta \sum_P \bigl(1 - \chi(U_P)\bigr)
= \beta \sum_P \bigl(1 - \tfrac{1}{2}\mathrm{Re}\,\mathrm{tr}\,U_P\bigr),
\label{eq:wilson_action}
\end{equation}
with the sum ranging over the plaquettes enumerated as in
Table~\ref{tab:plaquettes}.  The gauge coupling is $\beta = 4/g^2$ in
Wilson's convention: large~$\beta$ is weak coupling (continuum
perturbation), small~$\beta$ is strong coupling (strong-lattice
expansion).  The action is non-negative and vanishes iff every
plaquette holonomy is the identity (trivial vacuum).

\subsection{Gauge transformation}

A gauge transformation is a map $g: V \to \mathrm{SU}(2)$ acting on
link variables as
\begin{equation}
U_{u \to v} \; \longmapsto \; g_u U_{u \to v} g_v^{\dagger}.
\label{eq:gauge_transform}
\end{equation}
Under this action, the plaquette holonomy of any closed plaquette
transforms by conjugation, $U_P \mapsto g_{v_0} U_P g_{v_0}^{\dagger}$,
where $v_0$ is the base vertex.  The trace
$\mathrm{tr}\,U_P$ is therefore gauge-invariant, and so is the action
$S_W$.  Numerical verification on $n_{\max} = 3$ with random Haar link
variables and random Haar gauge transformations gives
$|S_W[U] - S_W[g U g^{-1}]| \leq 2.7 \times 10^{-15}$ (machine
precision; worst case over 100 seeded Haar trials, and
$1.8 \times 10^{-15}$ in the recorded reference run); the committed
tests assert $< 10^{-13}$ per trial --- see
\texttt{tests/test\_su2\_wilson\_gauge.py}, class
\texttt{TestGaugeInvariance} and
\texttt{test\_s3\_coulomb\_max\_n3\_endtoend}.

\subsection{Partition function}

The partition function is the Haar integral over
$\mathrm{SU}(2)^{|E|}$,
\begin{equation}
Z(\beta) = \int \prod_{e} dU_e \; \exp\bigl(-S_W[U]\bigr),
\label{eq:partition_function}
\end{equation}
with $dU_e$ the normalized bi-invariant Haar measure on
$\mathrm{SU}(2)$.  This integral is finite-dimensional and
well-defined for all $\beta \in \mathbb{R}$: no continuum
renormalization is required.  A direct implementation is via Monte
Carlo sampling of link variables weighted by $\exp(-S_W)$; an
analytical approach is via the SU(2) character expansion of
Section~\ref{sec:characters}.

% ===================================================================
\section{Character Expansion on SU(2)}
\label{sec:characters}
% ===================================================================

The SU(2) characters of the $(j+1)$-dimensional irreducible
representations are
\begin{equation}
\chi_j(U) = \frac{\sin\bigl((j+1)\theta\bigr)}{\sin\theta},
\qquad U = e^{i\theta\,\hat{n}\cdot\boldsymbol{\sigma}},
\label{eq:characters}
\end{equation}
with the orthogonality $\int dU\, \chi_j(U)\,\overline{\chi_k(U)} =
\delta_{jk}$.  The Wilson Boltzmann factor expands as a
Fegan--Menotti--Onofri series~\cite{fegan1983, menotti_onofri1981, creutz1983}
\begin{equation}
\exp\bigl(\beta \cdot \tfrac{1}{2}\mathrm{Re}\,\mathrm{tr}\,U\bigr)
= \sum_{j=0}^{\infty} d_j(\beta)\,\chi_j(U),
\qquad
d_j(\beta) = \frac{2(j+1)}{\beta}\,I_{j+1}(\beta),
\label{eq:character_expansion}
\end{equation}
where $I_n$ is the modified Bessel function of the first kind.  At
$\beta = 0$, $d_0 = 1$ and $d_j = 0$ for $j \geq 1$, so $Z(0) = 1$
(trivially, Haar measure is normalized).  At $\beta \to \infty$,
$d_j \to 1$ for all fixed $j$, and the Gaussian quadratic
approximation around the trivial vacuum dominates
(Section~\ref{sec:L1_kinetic}).

\paragraph*{Coefficient table at $n_{\max} = 3$.}  For the two primitive
4-cycles of Eq.~(\ref{eq:p_block_plaquettes}), the leading-order
plaquette-independent partition function
$Z_0(\beta) = e^{-\beta p}\bigl[d_0(\beta)\bigr]^p$ with $p = 2$
gives, together with the first excited coefficient $d_1(\beta) =
(4/\beta) I_2(\beta)$,

\begin{table}[h]
\centering
\caption{Character coefficients and leading-order partition function
on the $n_{\max} = 3$ Hopf graph.  The ratio $d_1/d_0$ exceeds unity
at $\beta \approx 3$, so the plaquette-independence approximation
loses accuracy for $\beta \gtrsim 5$.  Numerical data:
\texttt{debug/data/su2\_wilson\_gauge\_results.json}.}
\label{tab:char_coeffs}
\begin{tabular}{rrrrr}
\toprule
$\beta$ & $Z_0(\beta)$ & $d_0(\beta)$ & $d_1(\beta)$ & $d_1/d_0$ \\
\midrule
 0.1 & $8.21\times 10^{-1}$ & $1.0013$ & $0.0500$ & $0.050$ \\
 0.5 & $3.91\times 10^{-1}$ & $1.0316$ & $0.2552$ & $0.247$ \\
 1.0 & $1.73\times 10^{-1}$ & $1.1303$ & $0.5430$ & $0.480$ \\
 2.0 & $4.63\times 10^{-2}$ & $1.5906$ & $1.3779$ & $0.866$ \\
 5.0 & $4.30\times 10^{-3}$ & $9.7343$ & $14.0045$ & $1.439$ \\
10.0 & $5.88\times 10^{-4}$ & $534.20$ & $912.61$  & $1.708$ \\
\bottomrule
\end{tabular}
\end{table}

The entries of Table~\ref{tab:char_coeffs} set the scale at which
higher-rep corrections matter.  For $\beta \leq 2$, $d_1/d_0 < 1$ and
the leading-order approximation tracks Monte Carlo data
(Section~\ref{sec:results_MC}) to roughly $20\%$.  For $\beta \geq 5$,
higher-rep contributions compete with the trivial channel and the
leading-order truncation loses accuracy.

% ===================================================================
\section{Three Structural Results}
\label{sec:results}
% ===================================================================

\subsection{Result 1: $U(1)$ reduction to Paper 25 is exact}
\label{sec:u1_reduction}

Paper~25 defines the abelian $U(1)$ Wilson--Hodge action on the same
graph by assigning a phase $\theta_e \in [0, 2\pi)$ to each oriented
edge and setting
\begin{equation}
S_{U(1)}[\theta] = \beta \sum_P \bigl(1 - \cos \theta_P\bigr),
\qquad
\theta_P = \sum_{e \in P} \theta_e,
\label{eq:u1_action}
\end{equation}
with reverse-edge phases negated and node-local gauge shifts
$\theta_e \mapsto \theta_e + (\chi_{t(e)} - \chi_{s(e)})$.

\begin{proposition}[Maximal-torus reduction]
\label{prop:u1_reduction}
Let $\mathrm{T} \subset \mathrm{SU}(2)$ be the maximal torus,
$\mathrm{T} = \{\mathrm{diag}(e^{i\theta}, e^{-i\theta}) : \theta \in
[0, 2\pi)\}$.  Restricting the link variables of the
$\mathrm{SU}(2)$ Wilson action (\ref{eq:wilson_action}) to
$\mathrm{T}$ yields the $U(1)$ Wilson action of
Eq.~(\ref{eq:u1_action}):
\begin{equation}
S_W\bigl[\{U_e = \mathrm{diag}(e^{i\theta_e}, e^{-i\theta_e})\}\bigr]
= S_{U(1)}[\theta].
\label{eq:u1_reduction_claim}
\end{equation}
\end{proposition}

\begin{proof}
A product of diagonal $\mathrm{SU}(2)$ matrices is diagonal; by
induction on the plaquette length,
$U_P = \mathrm{diag}(e^{i\theta_P}, e^{-i\theta_P})$ where $\theta_P
= \sum_{e \in P} \theta_e$ (with reverse-edge phases counted
negatively, matching the convention $U_{v \to u} = U_{u \to v}^{\dagger}$
which for diagonal links negates $\theta_e$).  Then
$\tfrac{1}{2}\mathrm{Re}\,\mathrm{tr}\,U_P =
\tfrac{1}{2}(e^{i\theta_P} + e^{-i\theta_P}) = \cos \theta_P$, and
substitution into (\ref{eq:wilson_action}) yields (\ref{eq:u1_action}).
\qed
\end{proof}

The proposition is verified numerically in
\texttt{tests/test\_su2\_wilson\_gauge.py}, class
\texttt{TestU1Reduction}: for a triangle plaquette with random
diagonal phases $\{\theta_e\} = \{0.3, 0.7, -0.5\}$ and $\beta = 1.7$,
the two action evaluations agree to $|S_W - S_{U(1)}| < 10^{-12}$
(machine precision, 170 lines of test code).  A fuller
random-Haar-diagonal check at $\beta = 2$ on $n_{\max} = 3$ gives
$|S_W - S_{U(1)}| = 0$ exactly.

\begin{observation}[Paper 25 is the Cartan-subalgebra sector]
\label{obs:cartan_subalgebra}
Paper~25's abelian $U(1)$ Wilson--Hodge construction is structurally
the Cartan-subalgebra projection of the present non-abelian
$\mathrm{SU}(2)$ theory on the same graph.  The question ``is
Paper~25 abelian because the underlying geometry forces it to be, or
because abelian is one choice out of several?'' has a definite
answer: Paper~25 is the abelian sub-sector of a genuinely
non-abelian lattice gauge theory on the same graph.  This elevates
Paper~25's $U(1)$ construction from a freestanding abelian model to
the commutative-reduction sector of a natural non-abelian one.
\end{observation}

\begin{remark}[higher-rank generalization]
\label{rem:higher_rank}
On a rank-$r$ gauge group with maximal torus $\mathrm{T} = U(1)^r$,
the analog of Proposition~\ref{prop:u1_reduction} produces $r$
independent (but coupled, via a $\det = 1$ constraint when present)
$U(1)$ sectors.  The $\mathrm{SU}(3)$ case ($r = 2$) is computed in
Sprint~ST-SU3 (May~2026,
\texttt{debug/st\_su3\_wilson\_memo.md}~\S~3): restricting Wilson
$\mathrm{SU}(3)$ links on the Bargmann--Segal $S^5$ graph to the
maximal torus $T \cong U(1) \times U(1)$ produces a coupled
$U(1) \times U(1)$ Wilson action with character
$(1/3)\bigl(\cos a_P + \cos b_P + \cos(a_P + b_P)\bigr)$, verified
at machine precision over random Cartan trials at $N_\text{max} = 2,
3$ (max $|\Delta S| < 1.5 \times 10^{-11}$ at $7\,765$ plaquettes).
\end{remark}

The character expansion of Section~\ref{sec:characters} reflects the
reduction explicitly: the restriction of the $\mathrm{SU}(2)$
character $\chi_j$ to the maximal torus
$U = \mathrm{diag}(e^{i\theta}, e^{-i\theta})$ is
\begin{equation}
\chi_j\bigl(\mathrm{diag}(e^{i\theta}, e^{-i\theta})\bigr)
= \frac{\sin\bigl((j+1)\theta\bigr)}{\sin\theta}
= \sum_{k=0}^{j} e^{i(j - 2k)\theta},
\label{eq:character_restriction}
\end{equation}
a sum of $j+1$ $U(1)$ Fourier modes.  The Fegan--Menotti--Onofri
expansion (\ref{eq:character_expansion}) therefore restricts to the
classical Jacobi--Anger expansion of $\exp(\beta \cos \theta)$ on the
maximal torus---Paper~25's abelian expansion recovered term by term.

\subsection{Result 2: the weak-coupling kinetic term completes
$L_1 = B^{\top} B$ to the full Hodge-1 Laplacian}
\label{sec:L1_kinetic}

Paper~25~\S II--III introduces the edge Laplacian $L_1 = B^{\top} B$
as the ``discrete Hodge-1 Laplacian'' whose kernel dimension equals
$\beta_1$ and whose nonzero spectrum coincides with the nonzero
spectrum of $L_0$ (by singular-value decomposition of $B$), and notes
that on a bare 1-complex only the \emph{exact} part
$d_0 d_0^*$ of the manifold Hodge-1 operator
$\Delta_1 = d_0 d_0^* + d_1^* d_1$ is available --- the co-exact part
requires faces.  The present section identifies the \emph{kinetic
term} of the $\mathrm{SU}(2)$ Wilson action at weak coupling and
shows it is exactly that missing co-exact (curl) part, built on the
Ihara plaquettes as faces.

Write each link variable as an exponential of a
$\mathfrak{su}(2)$-valued gauge field,
\begin{equation}
U_e = \exp\bigl(i\, A_e^{a} \sigma_a / 2\bigr),
\qquad a = 1, 2, 3,
\label{eq:link_exponential}
\end{equation}
with $A_e^{a} \in \mathbb{R}$ small (weak-coupling regime).  To second
order,
\begin{equation}
U_e = I + \tfrac{i}{2} A_e^{a} \sigma_a - \tfrac{1}{8} A_e^{a} A_e^{b}
\sigma_a \sigma_b + O(A^3).
\label{eq:link_expansion}
\end{equation}
The plaquette holonomy $U_P$ at second order in $A$ is (using
$\sigma_a \sigma_b = \delta_{ab} I + i \epsilon_{abc} \sigma_c$ to
organize the Pauli products)
\begin{equation}
U_P = I + \tfrac{i}{2}\Bigl(\sum_{e \in P} A_e^{a}\Bigr) \sigma_a
- \tfrac{1}{8}\Bigl(\sum_{e \in P} A_e^{a}\Bigr)\Bigl(\sum_{e' \in P}
  A_{e'}^{a}\Bigr) I + O(A^3) + (\text{commutators}),
\label{eq:plaquette_expansion}
\end{equation}
where the ``commutators'' term is traceless and collects the
non-abelian corrections $[A_e, A_{e'}]$ that appear at third order
and higher.  Taking $\tfrac{1}{2} \mathrm{Re}\,\mathrm{tr}$ and
subtracting $1$ gives
\begin{equation}
1 - \tfrac{1}{2}\mathrm{Re}\,\mathrm{tr}\,U_P
= \tfrac{1}{8}\Bigl(\sum_{e \in P} A_e^{a}\Bigr)^2 + O(A^4)
\qquad (\text{sum over }a = 1, 2, 3 \text{ implicit}).
\label{eq:plaquette_quadratic}
\end{equation}
The $O(A^3)$ term vanishes because $\mathrm{tr}\,\sigma_a = 0$.  The
sum $\sum_{e \in P} A_e^{a}$ is exactly the signed sum of the gauge
field over the plaquette boundary, which is the action of the
discrete exterior derivative $d_1 = B^{\top}$ (acting on 1-cochains;
in fact the $d_1$ that maps 1-cochains to 2-cochains, but on a
triangle-free 1-complex we can read it as a signed sum over
plaquettes).  Therefore
\begin{equation}
S_W[U] = \frac{\beta}{8} \sum_P \Bigl(\sum_{e \in P} A_e^{a}\Bigr)^2
+ O(A^4)
= \frac{\beta}{8}\, A^{a\top} M A^{a} + O(A^4),
\label{eq:wilson_quadratic}
\end{equation}
where $M = B_2 B_2^{\top}$ is the plaquette-incidence bilinear form
on the edge space, built from the signed edge--face incidence matrix
$B_2$ whose columns are the (signed) indicator vectors of the Ihara
plaquettes.  In Hodge terms, $M = d_1^* d_1$ is the \emph{co-exact
(curl) part} of the Hodge-1 Laplacian of the 2-complex obtained by
attaching the plaquettes as faces.  It is \emph{not} proportional to
Paper~25's $L_1 = B^{\top} B = d_0 d_0^*$: the two quadratic forms
have complementary supports.  On the $n_{\max} = 3$ graph
($E = 13$, $\beta_1 = 2$) one computes $\operatorname{rank} M = 2 =
\beta_1$, $\operatorname{rank} L_1 = 11 = E - \beta_1$, and
$L_1 B_2 = 0$ exactly --- the plaquette columns span $\ker L_1$ (the
cycle space), which is precisely where the divergence form $L_1$
vanishes.  Their sum $L_1 + M$ has full rank $E$ and is the complete
discrete Hodge-1 Laplacian of the plaquette-filled 2-complex
(verified in \texttt{tests/test\_su2\_wilson\_gauge.py}, class
\texttt{TestWeakCouplingKinetic}).  The structural statement is:

\begin{proposition}[Wilson kinetic term = co-exact completion of $L_1$]
\label{prop:L1_kinetic}
The second-order expansion of the $\mathrm{SU}(2)$ Wilson action
(\ref{eq:wilson_action}) around the trivial vacuum is
\begin{equation}
S_W[U] = \tfrac{\beta}{8} \sum_a \langle A^{a}, B_2 B_2^{\top} A^{a}\rangle
+ O(A^4) + (\text{non-abelian})
\label{eq:L1_kinetic_form}
\end{equation}
where $B_2$ is the signed edge--plaquette incidence matrix (the
discrete $d_1$-boundary data supplied by the Ihara plaquette
enumeration).  The kinetic quadratic form $B_2 B_2^{\top}$ is the
co-exact (curl) part of the discrete Hodge-1 Laplacian of the
plaquette-filled 2-complex; it is supported exactly on
$\ker(L_1) = \ker(B^{\top} B)$, the $\beta_1$-dimensional cycle
space, and together with Paper~25's exact part it assembles the full
Hodge-1 operator $\Delta_1^{\text{disc}} = B^{\top} B +
B_2 B_2^{\top}$ (full rank at $n_{\max} = 3$).  The
$\mathfrak{su}(2)$ components $A^{a}$ ($a = 1, 2, 3$) are
kinetic-decoupled at quadratic order with the universal per-plaquette
per-component coefficient $1/8 = 1/(4 N_c)|_{N_c = 2}$; non-abelian
couplings first enter at order $A^4$.
\end{proposition}

\textbf{Provenance: MEASURED.}  The $1/8$ coefficient and the
quadratic form are verified numerically against exact
matrix-exponential holonomies, both on a triangle plaquette and on
the full $n_{\max} = 3$ graph (relative deviation $2.0 \times 10^{-3}$
at $\epsilon = 0.01$ shrinking to $7.3 \times 10^{-5}$ at
$\epsilon = 0.005$, the $O(A^4)$ remainder), in
\texttt{tests/test\_su2\_wilson\_gauge.py}::\texttt{TestWeakCoupling%
Kinetic} --- the $\mathrm{SU}(2)$ mirror of the $\mathrm{SU}(3)$
$1/12$ test in \texttt{tests/test\_su3\_wilson\_s5.py}.  An earlier
version of this proposition asserted that the kinetic form
``coincides with $L_1 = B^{\top} B$ up to a positive scalar
multiple''; that assertion was wrong (the two forms have
complementary supports, as the rank computation above shows) and is
withdrawn --- the correct statement is the completion statement,
which matches the face-structure reading already recorded in
Paper~25~\S II.

\begin{observation}[The Wilson kinetic sector completes the Hodge pair]
\label{obs:L1_structural}
At weak coupling the non-abelian Wilson theory contributes the
co-exact half of the Hodge-1 pair: the free (Gaussian) sector is the
curl form $B_2 B_2^{\top}$ on the cycle space, while Paper~25's
$L_1 = B^{\top} B$ is the exact half, whose role is combinatorial
(divergence / gradient sector) rather than plaquette-kinetic.  The
non-abelian self-couplings (cubic and quartic $A$-vertices) appear at
order $A^3$ and $A^4$ in the Wilson action expansion and involve the
structure constants of $\mathfrak{su}(2)$, namely $\epsilon^{abc}$;
they are parametrically suppressed at weak coupling.  The QED role of
the co-exact propagator $(B_2 B_2^{\top})^+$ --- recovering the
continuum self-energy structural zero that $L_1^+$ alone cannot ---
is established in Paper~28 (see Paper~25~\S II).
\end{observation}

\begin{remark}[universal $\mathrm{SU}(N)$ kinetic coefficient]
\label{rem:sun_coefficient}
With generators $T^a = \lambda^a/2$ normalized to
$\mathrm{tr}(T^a T^b) = \tfrac{1}{2} \delta^{ab}$, the second-order
expansion of the Wilson action
$S_W[U] = \beta \sum_P (1 - (1/N_c)\,\mathrm{Re}\,\mathrm{tr}\,U_P)$
on a general gauge group $\mathrm{SU}(N_c)$ produces a per-plaquette
per-component coefficient of $1/(4 N_c)$.  The $\mathrm{SU}(2)$ value
$1/8$ derived above is the $N_c = 2$ specialization, verified
numerically in \texttt{tests/test\_su2\_wilson\_gauge.py}::%
\texttt{TestWeakCouplingKinetic::test\_quadratic\_expansion\_numerically}.
The $\mathrm{SU}(3)$ value $1/12$ is verified independently in
Sprint~ST-SU3 (\texttt{debug/st\_su3\_wilson\_memo.md}~\S~4)
symbolically and numerically on the Bargmann--Segal $S^5$ graph
($120$ degrees of freedom at $N_\text{max} = 2$, relative error
$\sim O(\epsilon^2)$ consistent with $O(A^4)$ leading correction;
frozen test \texttt{tests/test\_su3\_wilson\_s5.py}::%
\texttt{TestWeakCouplingKinetic}).
The mechanism: more colors $\Rightarrow$ smaller per-component
coefficient because the $(1/N_c)\,\mathrm{Re}\,\mathrm{tr}$
normalization in the Wilson action divides by $N_c$ while the
trace-of-pair normalization holds $\mathrm{tr}(T^a T^a) = 1/2$ fixed.
The decoupling of components at quadratic order is a universal
$\mathrm{SU}(N_c)$ feature; non-abelian self-interactions enter at
$O(A^4)$ via the structure constants $f^{abc}$ in every case.
\end{remark}

\subsection{Result 3: Plaquette counts and Wilson-loop data}
\label{sec:results_MC}

Monte Carlo with Metropolis updates on the full
$\mathrm{SU}(2)^{|E|}$ configuration space provides direct access to
the Wilson-loop expectation value $\langle W_C \rangle = \bigl\langle
\tfrac{1}{2}\mathrm{Re}\,\mathrm{tr}\,U_C \bigr\rangle$.  We run at
$n_{\max} = 3$ (13 forward links), 2000 samples after 500 thermalization
steps, seed = 12, for a Wilson loop coinciding with one of the two
primitive 4-cycles in the $p$-block.

\begin{table}[h]
\centering
\caption{Wilson-loop expectation on the $n_{\max} = 3$ Coulomb graph,
loop = first $p$-block 4-cycle
[$(2,1,-1) \to (2,1,0) \to (3,1,0) \to (3,1,-1)$].  MC: Metropolis
sampling, 2000 samples, stderr reported; the seeded chain is
bit-deterministic and the four rows are pinned by the frozen test
\texttt{tests/test\_su2\_wilson\_gauge.py}::%
\texttt{TestMonteCarloWilsonLoop} (slow-marked).  Leading-order
character expansion:
$d_1(\beta) / d_0(\beta) = I_2(\beta) / I_1(\beta)$.  Numerical data:
\texttt{debug/data/su2\_wilson\_gauge\_results.json}.}
\label{tab:wilson_loop}
\begin{tabular}{rrrr}
\toprule
$\beta$ & $\langle W \rangle_{\mathrm{MC}}$ & $\pm\,$stderr & $\langle W \rangle_{\text{char}}$ \\
\midrule
0.5 & $0.126$ & $0.011$ & $0.124$ \\
1.0 & $0.158$ & $0.011$ & $0.240$ \\
2.0 & $0.332$ & $0.009$ & $0.433$ \\
5.0 & $0.681$ & $0.006$ & $0.719$ \\
\bottomrule
\end{tabular}
\end{table}

The two estimates agree at the $2$--$35\%$ level ($1.8\%$, $34\%$,
$23\%$, $5.3\%$ at $\beta = 0.5, 1, 2, 5$; loosest at intermediate
$\beta$; corrected 2026-07-03 from an earlier ``$\sim 20\%$''
characterization that understated the $\beta = 1$ row).  Discrepancy is structural: the leading-order character
expansion treats plaquettes $P_A$ and $P_B$ of
Eq.~(\ref{eq:p_block_plaquettes}) as independent, which they are not
(they share the undirected edge $\{(2,1,0), (3,1,0)\}$).  Monte Carlo
incorporates the shared-edge correlation.  Both estimates agree that
$\langle W \rangle$ is monotonic in $\beta$: strong coupling small
$\beta$ suppresses the trace, weak coupling large $\beta$ drives
$U_P \to I$ and $\langle W \rangle \to 1$.

\paragraph*{Confinement diagnostic on a single-plaquette-class graph.}
Confinement in Wilson's sense is the statement that $\langle W_C
\rangle$ obeys an \emph{area law} $\langle W_C \rangle \sim
e^{-\sigma \cdot \mathrm{Area}(C)}$ at strong coupling, crossing over
to a \emph{perimeter law} at weak coupling.  The diagnostic requires
a family of Wilson loops indexed by a geometric area.  On the
$n_{\max} = 3$ Coulomb graph, both primitive 4-cycles $P_A$ and
$P_B$ have the same length (4) and the same ``area'' (they span
$2 \times 2$ patches of the $m_l$-ladder / $n$-ladder lattice
within the $p$-block), so a Wilson-loop family is not available.
The Monte Carlo data do not (and cannot) show an area--perimeter
crossover at this size.  \emph{This is a structural feature of the
small graph, not a statement about the field theory}.  To see
confinement in the conventional sense, one needs a graph with a
continuous family of Wilson-loop sizes---i.e., $n_{\max} \geq 5$ or
finer discretizations of the continuum $S^3$.  Section~\ref{sec:open}
lists this as an immediate follow-up.

\paragraph*{Consistency with Paper 29's Ihara cycle counts.}  The
primitive 4-cycle count at $n_{\max} = 3$ ($N_4 = 2$) and the
higher-length counts ($N_6 = 1$, $N_8 = 1$) match the M\"obius-inverted
Ihara primitive-walk counts of Paper~29~\cite{paper29} up to the
orientation factor $N_{\text{M\"obius}}(L) = 2 \cdot N_{\text{unoriented}}(L)$
(each undirected cycle gives clockwise and counterclockwise primitive
walks; we keep one representative per unoriented class).  Paper~29's
plaquette-enumeration machinery is reused here without modification;
the Ihara cycles of one paper are the Wilson plaquettes of the other.

% ===================================================================
\section{Principle of Least Action: A Framework-Wide Synthesis}
\label{sec:action_principles}
% ===================================================================

The Wilson action of Eq.~(\ref{eq:wilson_action}) organizes the
gauge-sector dynamics of the GeoVac graph as a variational principle
on $\mathrm{SU}(2)^{|E|}$.  This is the fifth independent action
principle in the GeoVac framework, and it is useful to place it in
context.

Over the course of the project, variational or action structures have
emerged on five distinct data sectors of the Hopf graph.  Each is
tightly localized to one kind of structure and does not substitute
for any of the others.  We collect them in
Table~\ref{tab:action_principles}.

\begin{table}[t]
\centering
\small
\caption{Five action principles in the GeoVac framework.  Each
organizes a different sector of the graph data.  The Wilson-action
entry is the contribution of this paper.}
\label{tab:action_principles}
\begin{tabular}{>{\raggedright\arraybackslash}p{3.2cm}>{\raggedright\arraybackslash}p{3.1cm}>{\raggedright\arraybackslash}p{3.0cm}>{\raggedright\arraybackslash}p{3.2cm}}
\toprule
Principle & Data sector & Extremum principle & Reference \\
\midrule
Maximum entropy on packing & Nodes (candidate shell centers) & Maximize entropy under isotropy (Axiom 2) & Paper~0~\cite{paper0} \\
\addlinespace
Rayleigh--Ritz on graph Hamiltonian & Node wavefunctions (0-cochains) & Minimize $\langle \psi | H | \psi \rangle / \langle \psi | \psi \rangle$ & Papers~1, 7~\cite{paper7} \\
\addlinespace
Wilson lattice gauge action & Edge link variables (1-cochains, SU(2)-valued) & Minimize $S_W[U] = \beta \sum_P (1 - \tfrac{1}{2}\mathrm{Re}\,\mathrm{tr}\,U_P)$ & Paper~25~\cite{paper25} $U(1)$; present paper SU(2) \\
\addlinespace
Ihara zeta path integral & Primitive closed non-backtracking walks & $\log \zeta_G(s)^{-1} = \sum_{[C]} s^{\ell(C)}$ as weighted walk sum & Paper~29~\cite{paper29} \\
\addlinespace
Spectral action on Dirac spectrum & Spinor-bundle spectrum ($|\lambda_n|, g_n$) & $\mathrm{Tr}\,f\bigl(D / \Lambda\bigr)$ over modes & Paper~28~\cite{paper28} \\
\bottomrule
\end{tabular}
\end{table}

\paragraph*{What each sector is about.}
The node sector carries the scalar matter content (the wavefunction
amplitudes) and hosts two action principles: Paper~0's packing
axioms as a max-entropy statement (``maximize entropy under
isotropy''), which \emph{builds} the graph, and the Rayleigh--Ritz
principle on the built graph's Hamiltonian, which computes the
graph's spectrum.  The edge sector carries gauge content (Paper~25's
$U(1)$ phases, this paper's SU(2) link variables) and hosts the
Wilson action.  The walk sector aggregates edge data into
Ihara-cycle weights and hosts a discrete path integral
(log-partition function = walk sum).  The Dirac spectrum, finally,
is the specific data set on which the spectral action of
noncommutative-geometry type is computed; Paper~28's QED spectral
sums are literally this.

\paragraph*{The packing--Wilson duality.}
Two of the five principles stand out as structural duals.  Paper~0's
Axiom 2 (geometric indifference) is a max-entropy principle on node
data: among all candidate discrete packings compatible with Axiom 1
(binary distinguishability), the entropy-maximizing one is selected.
The Wilson action is a min-action principle on edge data: among all
SU(2) link configurations, those minimizing plaquette curvature are
selected in the continuum ($\beta \to \infty$) limit.  The two
principles live on complementary cochain degrees (0-form vs 1-form)
and extremize in opposite directions (max vs min), but they are
structurally compatible: Paper~0's max-entropy on nodes fixes
\emph{which} graph the Wilson action lives on, and the Wilson action
fixes \emph{how} gauge dynamics runs on it.  Nothing in the
mathematics forbids a joint variational principle on the full
(node, edge) data; we note this for future work but do not pursue it
here.

\paragraph*{What this section does and does not claim.}
The observation is descriptive: the GeoVac framework, as it currently
stands, presents five independent variational structures, each
localized to a different sector of the graph data.  We do not claim
that these five can be reduced to a single grand principle, nor that
one is more fundamental than the others.  Each is a tool matched to a
data type.  The utility of naming them together is that it makes the
project's full toolkit visible in one place: the framework is not
reducible to ``diagonalize a graph Laplacian,'' because its data has
more kinds of structure than nodes alone.

% ===================================================================
\section{Yang--Mills Adjacency: What Is and Is Not Claimed}
\label{sec:scope}
% ===================================================================

The present paper inherits the language of Wilson lattice gauge theory.
The rhetoric rule of CLAUDE.md~\S1.5 requires that we be explicit about
scope.  We enumerate what is done and what is not.

\paragraph*{What this paper establishes.}  A gauge-invariant, finite,
Haar-normalizable non-abelian lattice gauge theory on the finite
compact GeoVac Hopf graph at any $n_{\max}$, with (i) link variables
in $\mathrm{SU}(2)$, (ii) plaquettes as primitive closed
non-backtracking walks, (iii) Wilson action and partition function
well-defined without renormalization, (iv) exact $U(1)$ reduction to
Paper~25 in the maximal-torus limit, (v) the weak-coupling kinetic
term identified as the co-exact (curl) completion of
$L_1 = B^{\top} B$, with universal coefficient $1/8 = 1/(4N_c)$.
All 36 tests in \texttt{tests/test\_su2\_wilson\_gauge.py} pass,
including the numerical $U(1)$-reduction test, end-to-end gauge
invariance at machine precision (asserted $< 10^{-13}$), the kinetic
quadratic-form tests, and the seeded Monte Carlo Wilson-loop table.

\paragraph*{What this paper does NOT establish.}
\begin{itemize}
\item \emph{Yang--Mills mass gap on $\mathbb{R}^4$.}  The Clay
  Millennium Problem requires a continuum-field-theory construction
  of pure Yang--Mills on 4D Minkowski or Euclidean space, proof of a
  positive mass gap, and proof that no interacting scale-invariant
  continuum limit exists.  Our construction lives on a fixed finite
  compact graph; there is no continuum limit of $\mathbb{R}^4$ inside
  it.
\item \emph{Confinement in the conventional sense.}  As noted in
  Section~\ref{sec:results_MC}, the $n_{\max} = 3$ graph has a single
  plaquette-length class, so the area/perimeter-law diagnostic is
  degenerate.  Finite-size graphs with multiple plaquette classes
  (e.g., $n_{\max} \geq 5$) would be needed to even state the
  diagnostic.
\item \emph{Lorentzian signature or dynamics.}  This is a Euclidean
  lattice gauge theory on a spatial slice: no time direction, no
  Wick rotation, no $S$-matrix, no scattering.  Dynamics in the
  usual sense is absent.
\item \emph{Renormalization or continuum limit.}  The bare Wilson
  action is the whole story on a finite graph.  Whether the
  $n_{\max} \to \infty$ limit recovers continuum SU(2) Yang--Mills
  on the unit $S^3$, and in what sense, is an open question we do
  not address here.
\item \emph{No new derivation of $\alpha$.}  Paper~2's
  combination rule $K = \pi(B + F - \Delta)$ is not refined by the
  present construction.  Paper~25~\S IV read the three ingredients as
  Casimir-matter trace / Fock-Dirichlet gauge zeta / Dirac-boundary
  count; the non-abelian SU(2) Wilson structure does not produce a
  new $K$-like object or a non-abelian ``fine-structure constant''
  coupling, and we make no claim in that direction.
\end{itemize}

\paragraph*{Relationship to the Clay problem.}
The mathematical vocabulary is shared (lattice gauge theory,
non-abelian link variables, plaquette action, Wilson loops), but the
setting differs categorically.  The Clay problem is about continuum
field theory on a non-compact 4D spacetime; the present construction
is about discrete gauge theory on a compact 3D-graph spatial slice.
We use Wilson's formalism because it is the correct tool for
non-abelian lattice gauge theories on finite graphs; we do not
imply that the construction transports to $\mathbb{R}^4$.

% ===================================================================
\section{Open Questions}
\label{sec:open}
% ===================================================================

\subsection{Matter coupling: Dirac fermions on the Hopf graph}

The GeoVac framework's Tier-2 relativistic pipeline~\cite{paper14}
builds spin-ful composed qubit Hamiltonians in the
$(\kappa, m_j)$ Dirac basis on $S^3$.  These are matter sectors on
the same graph whose gauge sector we have built here.  Coupling the
two---Dirac fermion matter on nodes, $\mathrm{SU}(2)$ Wilson gauge on
edges---would be the first matter$+$gauge lattice theory in the
GeoVac project.  The construction is standard in principle (covariant
hopping amplitudes $\bar\psi_i U_{ij} \psi_j$), but on a
triangle-free simplicial 1-complex with Dirac spinors it would
require the Szmytkowski angular labels of Paper~14's Tier~2 to be
reconciled with the Wilson link-variable orientation convention.  A
clean demonstrator would be: does the spinor spectrum feel the
non-trivial $\mathrm{SU}(2)$ topology (i.e., is the Dirac operator
sensitive to the $\beta_1 = 2$ harmonic 1-forms at $n_{\max} = 3$ when
the gauge field is non-trivially topological)?

\subsection{Larger graphs and the confinement diagnostic}

At $n_{\max} \geq 5$, the Coulomb graph acquires multiple plaquette-length
classes, making a Wilson-loop family indexed by area available for
the first time.  The natural question: on the $n_{\max} = 5, 6$
Coulomb graphs, does the Wilson loop exhibit area-law behavior at
small $\beta$ and perimeter-law behavior at large $\beta$?  Running
this is computationally straightforward (the Metropolis update scales
as $|E|^2$ per sweep; $n_{\max} = 5$ has $|E| \sim 90$, manageable).
The result would distinguish a confining non-abelian theory from a
deconfined one on the GeoVac graph; Paper~25~\S VII's broader
questions about lattice-gauge structure would receive an immediate
sharpening.

\subsection{SU(2) analog of the Ihara zeta}

Paper~29 built the (scalar-valued) Ihara zeta of the GeoVac Hopf
graphs from primitive closed non-backtracking walk counts.  A natural
non-abelian generalization would weight each primitive walk $[C]$ by
a character $\chi_j(U_{[C]})$ of its $\mathrm{SU}(2)$ holonomy,
producing a character-valued zeta function.  Whether the resulting
object has a well-defined Ihara--Bass-type determinantal expression,
or whether it satisfies a graph-Riemann Hypothesis analog, is
speculative.  We record the question without pursuing it.

\subsection{SU(2) analog of Paper 2's $K = \pi(B + F - \Delta)$}

Paper~25~\S VII.A suggested that a non-abelian analog of the Paper~2
combination rule might exist on the $S^3$ graph.  The SU(2) Wilson
structure of the present paper does not by itself produce one; the
Wilson action is coupling-insensitive in a way that Paper~2's $K$ is
not.  A more promising route (if any) would couple the SU(2) Wilson
gauge to Dirac-spinor matter (first open question above) and examine
the combined spectral-gauge invariants.  We do not claim such a route
will succeed; we record it as a natural follow-up within the
Paper-2 observation-tier program.

\subsection{Plaquettes as the transverse photon sector}

The plaquettes enumerated in this paper define a boundary operator
$d_1\colon C^1 \to C^2$ whose transpose gives the co-exact Laplacian
$L_1^{\mathrm{co}} = d_1 d_1^\top$ on edges.  This is the
\emph{transverse} (physical photon) sector of the graph Hodge
decomposition, complementing the exact (longitudinal) sector
$L_1^{\mathrm{ex}} = B^\top B$ used in Paper~25 and the graph-native
QED of Paper~28.

Paper~28's pendant-edge theorem shows that the scalar propagator
$G = (B^\top B)^+$ breaks the continuum self-energy structural zero:
$\Sigma_{\mathrm{scalar}}(\mathrm{GS}) = 2(n_{\max}-1)/n_{\max}$.
Replacing $G$ with the transverse propagator
$G_T = (d_1 d_1^\top)^+$ recovers the structural zero \emph{exactly}:
$\Sigma_T(\mathrm{GS}) = 0$ at $n_{\max} = 3, 4$ (verified to machine
precision).  The mechanism is topological: the ground-state node is a
pendant vertex whose unique edge cannot participate in any closed walk,
so $d_1[e_0, :] = 0$ and hence $G_T[e_0, :] = 0$ identically.

This gives the plaquettes of this paper a concrete QED role beyond
their original gauge-theory motivation: they define the physical photon
sector that recovers continuum selection rules inaccessible from the
scalar (longitudinal) propagator alone.  One of the four
``vector-photon-required'' selection rules identified in Paper~28's
census is now ``plaquette-recoverable,'' refining the three-tier
partition to four tiers.

\subsection{Continuum limit}

The continuum limit of lattice gauge theory on refined Alon--Boppana
sequences of graphs converging to a target manifold is a technical
question that has been addressed in specific cases (hypercubic
lattices converging to $\mathbb{R}^4$, spherical lattices converging
to $S^3$).  Whether the $n_{\max} \to \infty$ sequence of GeoVac
graphs admits a well-posed continuum limit, and whether that limit is
SU(2) Yang--Mills on the unit $S^3$ in the Hodge--de Rham sense, is
open.  The relevant spectral data is available in
\texttt{geovac/hopf\_bundle.py}; a careful analysis has not been done.

\subsection{SU(3) Wilson sibling on $S^5$}
\label{sec:su3_sibling}

Paper~25~\S~VII.A also flagged as open whether a non-abelian extension
of the GeoVac Hopf $U(1)$ Wilson--Hodge structure transfers to the
Bargmann--Segal $S^5$ graph (Paper~24).  Sprint~5 Track~S5
(\texttt{debug/s5\_gauge\_structure\_memo.md}) returned a clean
negative for the \emph{adjacency-preserving} SU(3) action on the
$(N, 0)$ representation tower: shell-shifting transitions are
Clebsch--Gordan intertwiners between distinct SU(3) irreps, not group
elements, so an SU(3) action on the $(N, 0)$ tower fails the Wilson
``fixed group on every link'' requirement.

Sprint~ST-SU3 (May~2026,
\texttt{debug/st\_su3\_wilson\_memo.md}, with module
\texttt{geovac/su3\_wilson\_s5.py} and tests
\texttt{tests/test\_su3\_wilson\_s5.py} 49/49 pass, two slow-marked) tested a
structurally different question: \emph{Wilson lattice gauge theory}
with link variables $U_e \in \mathrm{SU}(3)$ as external 3$\times$3
matrices in the fundamental representation, independent of the shell
index $N$.  The construction parallels the present paper's
$\mathrm{SU}(2)$ on $S^3$.  Three structural results mirror those of
the present paper:

\begin{itemize}
\setlength\itemsep{0pt}
\item \textbf{Cartan reduction.}  Restricting Wilson SU(3) links to
the maximal torus $T \cong U(1)\times U(1)$ produces a coupled
$U(1)\times U(1)$ Wilson action with character
$(1/3)\bigl(\cos a_P + \cos b_P + \cos(a_P+b_P)\bigr)$ (the third
phase forced by $\det = 1$).  Verified at machine precision over
random Cartan trials at $N_\text{max} = 2, 3$
(max $|\Delta S| < 1.5\times 10^{-11}$ at $7\,765$ plaquettes).
The rank-2 generalization of Proposition~\ref{prop:u1_reduction} per
Remark~\ref{rem:higher_rank}.
\item \textbf{Kinetic coefficient.}  The weak-coupling expansion gives
$1/12 = 1/(4 \cdot 3)$ per plaquette per generator component, the
$N_c = 3$ specialization of Remark~\ref{rem:sun_coefficient}'s
universal $1/(4 N_c)$ formula.  Eight $\mathfrak{su}(3)$ components
are kinetic-decoupled at quadratic order; non-abelian self-couplings
enter at $O(A^4)$ via the structure constants $f^{abc}$.
\item \textbf{Plaquette tabulation.}  $V, E, \beta_1$ at
$N_\text{max} = 2, 3, 4$ are $(10, 15, 6)$, $(20, 42, 23)$,
$(35, 90, 56)$, with primitive 4-cycle counts $N_4 = 15, 88, 272$
--- roughly an order of magnitude denser than the $S^3$ graph at
matched cutoff.  Wilson loops at $N_\text{max} = 2$ with local SU(3)
updates give $\langle W \rangle = 0.05, 0.10, 0.26, 0.65, 0.89$ at
$\beta = 0.5, 1, 2, 5, 10$, systematically lower than the present
paper's SU(2) values at small $\beta$ (consistent with the standard
$N_c$-color randomization of the trace).
\end{itemize}

The Sprint~5 Track~S5 Clebsch--Gordan obstruction is \emph{bypassed}
at the gauge level (Wilson links live in the fundamental SU(3) rep,
not in the $(N, 0)$ shell rep) but \emph{rediscovered} at the
matter-coupling level: coupling Wilson SU(3) to natural shell matter
(in the $(N, 0)$ irreps of growing dimension) requires a CG
intertwiner from $\textbf{3}\otimes (N, 0)$ to $(N+1, 0)$ that is
exactly the Sprint~5 obstruction.  Wilson SU(3) on Bargmann is
therefore a self-consistent \emph{pure-gauge} theory; the natural
matter sector cannot be coupled in the strict Wilson sense.

The structural reading: $S^3 = \mathrm{SU}(2)$ makes the
present paper's case special in a way that does not transfer to the
$\mathrm{SU}(3)$ on $S^5$ analogue.  On $S^3$, matter (spin-$1/2$
fermions, Paper~14 Tier~2) lives in the same representation as the
link variables, so the gauge--matter coupling is natural; on $S^5$,
the matter representations $(N, 0)$ have dimensions that grow with
$N$, so a uniform-link-rep coupling is impossible.  This sharpens
Paper~24's Coulomb/HO asymmetry from three layers
(spectrum-computing $L_0$, calibration $\pi$, Wilson gauge with
matter coupling) plus the universal Wilson--Hodge vocabulary, to
\emph{four layers}: Wilson gauge alone is universal; only Wilson
gauge \emph{with natural matter coupling} is Coulomb-specific.
Sprint~ST-SU3 contributes the third layer on the gauge-with-matter
side.  The reading of Paper~24 \S~V is updated accordingly.

% ===================================================================
\section{Conclusion}
\label{sec:conclusion}
% ===================================================================

Sprint 4 Track RH-Q constructed SU(2) Wilson lattice gauge theory on
the GeoVac Hopf graph (module \texttt{geovac/su2\_wilson\_gauge.py},
36 tests passing, 2026-07-04 count).  The present paper formalizes that construction
and places it in context relative to Paper~25's abelian $U(1)$
predecessor.

Three structural results summarize the content:
(1)~the maximal-torus reduction of the $\mathrm{SU}(2)$ Wilson action
recovers Paper~25's $U(1)$ Wilson--Hodge action exactly, so
Paper~25 is the Cartan-subalgebra sector of a non-abelian theory;
(2)~the weak-coupling kinetic term of the non-abelian Wilson action
is the co-exact (curl) quadratic form on the Ihara plaquettes ---
coefficient $1/8 = 1/(4N_c)$ per plaquette per
$\mathfrak{su}(2)$ component --- which completes the edge Laplacian
$L_1 = B^{\top} B$ that Paper~25 introduced to the full discrete
Hodge-1 Laplacian, with non-abelian self-interactions entering only
at quartic order in the gauge field;
(3)~plaquette enumeration on the $S^3$ Coulomb graph matches
Paper~29's Ihara cycle counts, and direct Monte Carlo on $n_{\max} = 3$
gives monotonic $\beta$-dependence of the Wilson loop expectation
(confirming the character-expansion leading order at the $\sim 20\%$
level, as expected from the shared-edge correction).

We close with a framework observation.  The $S^3 = \mathrm{SU}(2)$
coincidence is structurally overdetermined: $S^3$ appears in the
framework as Fock's projection of the Coulomb problem (Paper~7), as
the underlying manifold of the $(n, l, m_l)$ Hopf graph (Paper~14),
as the spin-1/2 carrier of the Dirac sector (Paper~14 Tier-2), and
now as the gauge group of the non-abelian Wilson construction.  Four
independent structural roles on the same 3-manifold, each needing
only the simplest compact simply-connected non-abelian Lie group of
rank~1.  We do not read this as evidence that $S^3$ is fundamental in
any ontological sense (CLAUDE.md~\S1.5 rhetoric); we do observe that
$S^3$ is, in a specific technical sense, the minimal object that
supports all four roles simultaneously, and we record the convergence
as a fact.

The project now contains five action principles---Paper~0's
max-entropy on nodes, Rayleigh--Ritz on the graph Hamiltonian,
Paper~25's and this paper's Wilson action on edges, Paper~29's Ihara
zeta path integral on walks, and Paper~28's spectral action on the
Dirac spectrum---each matched to a distinct sector of the graph
data.  None reduces to any of the others.  The Wilson action has,
with this paper, a genuine non-abelian sibling to its $U(1)$
predecessor; the rest of the table is left as is.

% ===================================================================
% References
% ===================================================================

\begin{thebibliography}{99}

\bibitem{wilson1974}
K.~G.~Wilson, ``Confinement of quarks,''
\textit{Phys.~Rev.~D}\ \textbf{10}, 2445--2459 (1974).

\bibitem{creutz1983}
M.~Creutz, \textit{Quarks, Gluons and Lattices},
Cambridge University Press, Cambridge (1983).

\bibitem{menotti_onofri1981}
P.~Menotti and E.~Onofri, ``The action of SU(N) lattice gauge
theory in terms of the heat kernel on the group manifold,''
\textit{Nucl.~Phys.~B}\ \textbf{190}, 288--300 (1981).

\bibitem{fegan1983}
H.~D.~Fegan, ``The heat equation on a compact Lie group,''
\textit{Trans.~Amer.~Math.~Soc.}\ \textbf{246}, 339--357 (1978).

\bibitem{lim2020}
L.-H.~Lim, ``Hodge Laplacians on graphs,''
\textit{SIAM Review}\ \textbf{62}, 685--715 (2020).

\bibitem{schaub2020}
M.~T.~Schaub, A.~R.~Benson, P.~Horn, G.~Lippner, and A.~Jadbabaie,
``Random walks on simplicial complexes and the normalized Hodge 1-Laplacian,''
\textit{SIAM Review}\ \textbf{62}, 353--391 (2020).

\bibitem{marcolli_vs2014}
M.~Marcolli and W.~D.~van Suijlekom,
``Gauge networks in noncommutative geometry,''
\textit{J.~Geom.~Phys.}\ \textbf{75}, 71--91 (2014),
arXiv:1301.3480.

\bibitem{perez_sanchez2025}
C.~I.~Perez-Sanchez,
``Comment on `Gauge networks in noncommutative geometry','',
arXiv:2508.17338 (2025).

\bibitem{garofalo2023_su2_s3}
M.~Garofalo, T.~Hartung, T.~Jakobs, K.~Jansen, J.~Ostmeyer, D.~Rolfes,
S.~Romiti, and C.~Urbach,
``Testing the SU(2) lattice Hamiltonian built from $S_3$ partitionings,''
arXiv:2311.15926 (2023).

\bibitem{jakobs2025_su2_dynamics}
T.~Jakobs, M.~Garofalo, T.~Hartung, K.~Jansen, J.~Ostmeyer, S.~Romiti,
and C.~Urbach,
``Dynamics in Hamiltonian lattice gauge theory: Approaching the
continuum limit with partitionings of SU(2),''
arXiv:2503.03397 (2025).

\bibitem{paper0}
GeoVac Project, ``Angular Momentum as Information Geometry: Isotropic
Packing and the Universal Angular Cross-Section,'' Paper~0 (2025).

\bibitem{paper2}
GeoVac Project, ``The Fine Structure Constant as a Spectral Coincidence
on the Hopf Fibration,'' Paper~2 (2025).

\bibitem{paper7}
GeoVac Project, ``The Dimensionless Vacuum: Recovering the Schr\"odinger
Equation from Scale-Invariant Graph Topology,'' Paper~7 (2025).

\bibitem{paper14}
GeoVac Project, ``Structurally Sparse Qubit Hamiltonians from Spectral
Graph Theory,'' Paper~14 (2025).

\bibitem{paper18}
GeoVac Project, ``Spectral-Geometric Exchange Constants: Projecting
Discrete Spectra onto Continuous Manifolds,'' Paper~18 (2026).

\bibitem{paper22}
GeoVac Project, ``Potential-Independent Angular Sparsity for Qubit
Hamiltonians in Spherical Fermion Systems,'' Paper~22 (2026).

\bibitem{paper24}
GeoVac Project, ``The Bargmann--Segal Lattice: $\pi$-Free
Discretization of the Harmonic Oscillator on $S^5$,'' Paper~24
(2026).

\bibitem{paper25}
GeoVac Project, ``The Hopf Graph as a Lattice Gauge Structure: A
Framework Observation,'' Paper~25 (2026).

\bibitem{paper28}
GeoVac Project, ``QED on $S^3$: Spectral Geometry of Perturbative
Transcendentals,'' Paper~28 (2026).

\bibitem{paper29}
GeoVac Project, ``Finite-Size Graph-RH for the GeoVac Hopf Graphs: Ihara Zeta, Bound-Crossing, and a Scope Boundary on Selberg-on-Hydrogen,'' Paper~29
(2026).

\bibitem{rh_q_memo}
GeoVac Project, ``Track RH-Q: SU(2) Wilson Lattice Gauge Theory on
the $S^3 = \mathrm{SU}(2)$ GeoVac Hopf Graph,''
\texttt{debug/su2\_wilson\_gauge\_memo.md}, April 2026.

\end{thebibliography}

\end{document}
